\documentclass[aspectratio=169,handout,notes,draft]{beamer}

\usetheme{Madrid}
\usecolortheme{beaver}
\usefonttheme{professionalfonts}
\useinnertheme{rounded}
\useoutertheme{infolines}

\setbeamertemplate{navigation symbols}{}
\setbeamercolor{structure}{fg=blue!60!black}
\setbeamerfont{frametitle}{size=\Large,series=\bfseries}
\setbeameroption{show notes}

\usepackage{graphicx}
\usepackage{amsmath}
\usepackage{booktabs}
\usepackage{tikz}
\usepackage{listings}
\usepackage[style=authoryear]{biblatex}
\addbibresource{refs.bib}

\newtheorem{claim}{Claim}

\title[Short title]{A talk about slides}
\subtitle{Frames, columns, blocks and overlays}
\author[A. Writer]{Ada Writer\inst{1} \and Ben Reader\inst{2}}
\institute[Univ.]{\inst{1}University of Somewhere \and \inst{2}Institute of Elsewhere}
\date{\today}
\titlegraphic{\includegraphics[height=1cm]{images/logo}}

\begin{document}

\begin{frame}{Blocks with and without titles}
\begin{block}{Block with title}
Plain block text.
\end{block}
\begin{block}{}
Block with an empty title.
\end{block}
\begin{alertblock}{}
Alert block with an empty title.
\end{alertblock}
\begin{exampleblock}{Example with title}
Example text.
\end{exampleblock}
\begin{block}<2->{Block with overlay}
Overlay block.
\end{block}
\begin{exampleblock}<3>{Example overlay}
Example overlay text.
\end{exampleblock}
\end{frame}

\begin{frame}<2->{Frame overlay and title}
Frame with an overlay specification and a title argument.
\end{frame}

\begin{frame}{Overlay theorem}
\begin{theorem}<2->[Named]
A theorem with an overlay and a name.
\end{theorem}
\begin{lemma}<3>
A lemma on slide three.
\end{lemma}
\end{frame}

\end{document}
